% =========================================================================
% WORKING DRAFT (restructured 2026-11-22 around measured results).
% DO NOT SUBMIT until:
%   (1) every [TO FILL] is completed with real experiments (baselines),
%   (2) every citation is re-verified against arXiv/DBLP/the PDF,
%   (3) a fresh literature sweep confirms the remaining "first" statements,
%   (4) the text is rewritten in the author's own voice.
% All result tables are \input from generated_tables.tex, written by
% make_tables.py from the saved result files. Never edit those numbers by hand.
% See INTEGRITY_CHECKLIST.md.
% =========================================================================
\documentclass[10pt,twocolumn]{article}
\usepackage[margin=0.75in,columnsep=0.25in]{geometry}
\usepackage{mathptmx}
\usepackage{amsmath,amssymb}
\usepackage{graphicx}
\usepackage{xcolor}
\usepackage[colorlinks=true,linkcolor=blue,citecolor=blue,urlcolor=blue]{hyperref}

\newcommand{\tofill}[1]{\textcolor{red}{\textbf{[TO FILL: #1]}}}

\title{\textbf{Bounding False Accusations in AI-Generated Image Detection:\\
an Audited Guarantee, an Honest Benchmark, and a Negative Result on Lens Physics}}
\author{Jitendra Jha\\
\small Independent Researcher\\
\small \texttt{github.com/JitendraJha98/fakeradar}}
\date{Working draft --- \today}

\begin{document}
\maketitle

\begin{abstract}
Detectors of AI-generated images are usually reported by ranking metrics on
benchmarks that resemble their training data, while the decision that matters
in practice --- flagging a real photograph as fake --- is left to a fixed 0.5
threshold. We study three questions. First, can single-image lens physics
serve as generator-independent evidence? We test whether defocus blur follows
the thin-lens law against monocular depth, with acceptance criteria fixed in
advance, and find that it cannot be measured reliably: real wide-aperture
photographs fit the law worse than images with deliberately broken blur, with
both a classical and a learned blur estimator. Second, can the rate of false
accusations be bounded without any assumption about generators? Using
split-conformal thresholds calibrated on real photographs only, stratified by
sharing platform, the realised false-accusation rate on 11 smartphones never
used for calibration stays at the target (4.7--5.9\% at a 5\% target,
0.8--1.3\% at 1\%) across camera originals and four platforms, and on a fifth
platform held out of calibration. Third, what does a small detector achieve
when the JPEG/PNG shortcut of GenImage is closed? A 12M-parameter
gradient-and-frequency detector trained on Stable Diffusion v1.4 reaches a
mean AUROC of 0.937 (three seeds; Table~\ref{tab:main}) over
eight generators, near-perfect on its training family and clearly weaker on
GAN- and ADM-based generators. Code, weights and every raw
score are released.
\end{abstract}

\section{Introduction}
Photorealistic generators have made fabricated imagery cheap. Detectors fail
in deployment along two well-documented axes: they degrade on generators not
seen in training~\cite{wang2020cnn,zhu2023genimage}, and on images that have
passed through compression and resizing~\cite{gragnaniello2021gan,corvi2023detection}.
A third failure receives less attention: the probability a detector outputs
is calibrated, if at all, on data resembling its training set, so a fixed
threshold carries no guarantee about how often a \emph{real} photograph is
flagged once the data shift. For the people affected, that false-accusation
rate is the number that matters.

This paper reports three results, each obtained under acceptance criteria
written down before any score was computed.
\begin{enumerate}\itemsep2pt
  \item \textbf{A negative result on lens-law evidence} (\S\ref{sec:lens}).
  A photograph is formed through one lens, so blur size is proportional to
  $|D - D_f|$ for inverse depth $D$. We test whether this law can be measured
  in single images well enough to expose generated images. It cannot:
  Gate~M fails with both a classical and a learned blur estimator.
  \item \textbf{An audited false-accusation guarantee} (\S\ref{sec:fpr}).
  Split-conformal thresholds from real photographs only, stratified by
  sharing platform, keep the realised rate at the target on phones never used
  for calibration (Gate~B passes, 10 of 10 checks).
  \item \textbf{An honest small-detector benchmark} (\S\ref{sec:bench}).
  Closing GenImage's compression shortcut~\cite{grommelt2024fake} and
  reporting three seeds, raw scores and an ablation, we measure what a
  12M-parameter CPU-class detector achieves.
\end{enumerate}

\section{Related Work}
\paragraph{Low-level and semantic detectors.}
High-pass residuals~\cite{fridrich2012rich}, gradients of a pretrained
network~\cite{tan2023lgrad}, up-sampling residuals~\cite{tan2024npr} and
spectral statistics~\cite{durall2020watch,frank2020leveraging} target
generation artifacts; frozen foundation-model features generalise better
across generators~\cite{ojha2023universal,cozzolino2024raising,liu2024fatformer}.
Hybrids that fuse both are established~\cite{yan2025sanity,cheng2025cospy}; our
detector is such a hybrid and we claim no novelty for its architecture.

\paragraph{Physical and geometric cues.}
Generated images violate projective geometry and shadow
constraints~\cite{sarkar2024shadows}. Defocus maps have been used as learned
features for deepfake detection~\cite{jeon2025defocus}, without depth or a
lens-law test. We test the thin-lens law directly.

\paragraph{Dataset biases.}
GenImage stores real photographs as JPEG and generated images as PNG, and
detectors learn compression as a shortcut~\cite{grommelt2024fake}.

\paragraph{Calibration and guarantees.}
Post-hoc calibration of synthetic-image detectors has been
studied~\cite{guo2017calibration}, and quality-aware calibration over
near-duplicate copies was proposed recently~\cite{guillaro2026quad}.
Conformal prediction gives distribution-free error control under
exchangeability~\cite{vovk2005algorithmic,angelopoulos2023gentle} and has been
used for watermark detection and machine-generated text; we did not find it
applied to unwatermarked AI-image detection with an audit across devices and
platforms \tofill{re-run the literature sweep before submission}.

\section{Lens-Law Evidence: A Negative Result}
\label{sec:lens}
\paragraph{Test.} For each image we estimate blur per $32{\times}32$ cell and
inverse depth with Depth Anything V2 Small~\cite{yang2024depthv2}, fit
$b(D) = \tfrac12\log(c_0^2 + k^2(D-D_f)^2)$ to log blur, and use the robust
fit residual as the score; images without visible depth of field abstain.

\paragraph{Gate M.} Before scoring any image we fixed five checks: real
aperture pairs must fit the V (M1); a wide-aperture shot must show more blur
spread than its narrow twin (M2); controlled renders of real photographs with
deliberately broken blur (step, image-row, missing foreground) must be
separable from lawful renders (M3); at a threshold flagging 5\% of real
wide-aperture photographs, at least half the broken renders must be flagged
(M4); and the separation must survive 0.5$\times$ downscaling with JPEG 70
(M5). Data: 349 iDFD indoor pairs with measured depth, 631 Unsplash
photographs shot at f/2.8 or wider, and renders from 248 photographs at f/8 or
narrower, all from the development split.

\paragraph{Estimators.} A classical multi-scale edge gradient-ratio estimator
after Zhuo and Sim~\cite{zhuo2011defocus}, and a learned estimator (0.92M
parameters) trained only on real sharp photographs with synthetic defocus
applied in linear light, which predicts log blur and its uncertainty per cell.

% Generated by paper/make_tables.py -- do not edit by hand.
\begin{table}[t]\centering\small
\caption{Gate M (development split). Both blur estimators fail; the test split was never scored.}
\label{tab:gatem}
\resizebox{\columnwidth}{!}{\begin{tabular}{l l l l l}\hline
Check & What & Needed & Edge ratio (v0) & Learned (v1)\\\hline
M1 & real pairs fit the V & $\ge$80\% / 70\% & 2.7\% / 2.3\% & 2.5\% / 2.2\%\\
M2 & wide shot has more blur spread & $\ge$90\% & 29.2\% & 99.4\%\\
M3 & AUROC, lawful vs.\ broken & $\ge$0.85 each & 0.76 / 0.69 / 0.65 & 0.85 / 0.78 / 0.66\\
M4 & broken / lawful flagged at 5\% & $\ge$50\% / $\le$10\% & 6.3\% / 4.5\% & 4.9\% / 1.4\%\\
M5 & AUROC after 0.5$\times$+JPEG70 & drop $\le$0.10 & 0.59 / 0.58 / 0.62 & 0.61 / 0.59 / 0.65\\
\hline\end{tabular}}
\end{table}


\paragraph{Result.} Gate~M fails with both estimators
(Table~\ref{tab:gatem}). The learned estimator fixes M2 and makes sharp shots
abstain (4\% applicable versus 83\%), but real wide-aperture photographs fit
the lens law \emph{worse} than the broken renders (median residual 0.40
versus 0.27--0.39), so a threshold set on real photographs flags almost
nothing; transport makes it worse. A stronger depth model of the same family
and a metric depth model did not help. On iDFD's indoor pairs defocus barely
varies with depth, so M1 was also run on 310 DPDD scenes~\cite{abuolaim2020dpdd},
where it does: 59--63\% of wide-aperture photographs fit the V (needed 80\% /
70\%) and 96--99\% show more blur spread than their stopped-down twin. The
law is therefore measurable in most strongly defocused photographs, but not
reliably enough per image to serve as evidence. Whether generators obey it remains open; a
stereo or dual-pixel measurement would be needed.

\section{An Audited False-Accusation Guarantee}
\label{sec:fpr}
\paragraph{Method.} Given scores $s_1,\dots,s_n$ of real calibration photos
in stratum $k$ (here: camera original or the sharing platform), the
threshold $t_k$ is the $\lceil (n+1)(1-\alpha)\rceil$-th smallest score; a
photo is flagged when its score exceeds $t_k$. If a new real photo is
exchangeable with the calibration photos of its stratum, it is flagged with
probability at most $\alpha$~\cite{vovk2005algorithmic}. No generated image
enters, so the bound cannot be broken by a new generator, only by a shift in
the real photographs; strata too small for level $\alpha$ use a pooled
threshold.

\paragraph{Audit (Gate B).} Detector: Community Forensics~\cite{park2025community}
(public MIT weights). Data: FODB~\cite{hadwiger2021forchheim}, 143 scenes
from 27 smartphones, each as camera original and as shared through Facebook,
Instagram, Telegram, Twitter and WhatsApp; 23{,}105 photos scored. Phones were
split 16/11 into calibration and audit by a fixed hash; Telegram was never
used for calibration. The pass rule was fixed in advance: a check passes if
the realised rate is at most $1.5\times$ the target or its 95\% phone-bootstrap
interval reaches the target; the gate passes if at most one of ten checks
misses and none exceeds $2\times$.

% Generated by paper/make_tables.py -- do not edit by hand.
\begin{table}[t]\centering\small
\caption{Gate B: share of real photos flagged on 11 audit phones never used for calibration (23{,}105 FODB photos). Passed: 0 of 10 checks missed.}
\label{tab:gateb}
\begin{tabular}{l c c}\hline
Version & target 5\% & target 1\%\\\hline
camera original & 5.5\% & 0.9\%\\
Facebook & 5.9\% & 1.3\%\\
Instagram & 5.5\% & 0.9\%\\
Twitter & 4.7\% & 1.0\%\\
WhatsApp & 5.5\% & 0.8\%\\
Telegram (held out) & 3.6\% & 0.9\%\\
\hline\end{tabular}
\end{table}


\paragraph{Result.} Gate~B passes with no misses (Table~\ref{tab:gateb}).
Stratification matters: a single pooled threshold flags 10.0\% of camera
originals at a 5\% target, because platform re-encoding shifts real scores.
The detector's own 0.5 cut flags almost no real photographs (0--0.4\%) but
also misses more generated images after transport (Table~\ref{tab:power}).
Limits: one detector; 2020-era phones and platforms; generated images for the
detection-rate step come from one generator (SDXL), likely within the
detector's training distribution, with simulated transport.

% Generated by paper/make_tables.py -- do not edit by hand.
\begin{table}[t]\centering\small
\caption{Share of 200 SDXL images flagged by the reference detector: its own 0.5 cut vs.\ thresholds calibrated on real photos (range over the five platform versions).}
\label{tab:power}
\begin{tabular}{l c c c}\hline
Condition & 0.5 cut & calibrated 5\% & calibrated 1\%\\\hline
as generated & 96.5\% & 99.0\%--100.0\% & 98.0\%--99.0\%\\
0.75$\times$ + JPEG 85 & 90.5\% & 99.5\% & 94.0\%--99.0\%\\
0.5$\times$ + JPEG 70 & 81.5\% & 96.5\%--98.0\% & 90.5\%--94.5\%\\
\hline\end{tabular}
\end{table}


\section{A Small Detector, Benchmarked Honestly}
\label{sec:bench}
\paragraph{Model.} A gradient-field branch (SRM residuals, an up-sampling
residual~\cite{tan2024npr}, gradient magnitude and structure-tensor
coherence of the luma gradient) on a ResNet-18~\cite{he2016resnet}, a radial
power-spectrum branch, a learned global softmax gate, auxiliary per-branch
heads and temperature calibration; native-resolution crops, never resized;
about 12M parameters.

\paragraph{Protocol.} Train on GenImage~\cite{zhu2023genimage} Stable
Diffusion v1.4 only: 95{,}000 images sampled from 324{,}000, 5\% more held out
for model selection and calibration, 5 epochs on one T4 GPU, with
transport-chain augmentation. Test on GenImage's \texttt{val} split of all
eight generators: 1{,}000 generated images per generator against 2{,}000
pooled real images. \emph{Equalized} runs JPEG-encode every non-JPEG image
once at quality 96 in training and testing, removing the compression
shortcut.

% Generated by paper/make_tables.py -- do not edit by hand.
\begin{table}[t]\centering\small
\caption{Cross-generator AUROC, trained on SD v1.4 only (mean$\pm$std over 2 seeds). ``Equalized'': non-JPEG images JPEG-encoded once at quality 96, closing the JPEG/PNG shortcut.}
\label{tab:main}
\begin{tabular}{l c c c}\hline
Test generator & equalized & as stored & w/o coherence\\\hline
SD v1.4 & 1.000$\pm$0.000 & 0.998$\pm$0.000 & 1.000$\pm$0.000\\
SD v1.5 & 0.999$\pm$0.000 & 0.998$\pm$0.000 & 0.999$\pm$0.000\\
Wukong & 0.995$\pm$0.000 & 0.980$\pm$0.006 & 0.996$\pm$0.000\\
VQDM & 0.922$\pm$0.001 & 0.844$\pm$0.021 & 0.918$\pm$0.005\\
Midjourney & 0.920$\pm$0.008 & 0.925$\pm$0.012 & 0.929$\pm$0.005\\
GLIDE & 0.893$\pm$0.000 & 0.844$\pm$0.004 & 0.886$\pm$0.006\\
BigGAN & 0.897$\pm$0.010 & 0.830$\pm$0.001 & 0.878$\pm$0.002\\
ADM & 0.862$\pm$0.031 & 0.793$\pm$0.029 & 0.836$\pm$0.027\\
\textbf{Mean of 8} & 0.936$\pm$0.001 & 0.902$\pm$0.001 & 0.930$\pm$0.006\\
\hline\end{tabular}
\end{table}

\begin{table}[t]\centering\small
\caption{Transport robustness: AUROC on 500 real + 500 generated test images (equalized; 2 seeds).}
\label{tab:robust}
\begin{tabular}{l c c}\hline
Condition & full model & w/o coherence\\\hline
clean & 0.935$\pm$0.001 & 0.926$\pm$0.003\\
JPEG 90 & 0.886$\pm$0.008 & 0.878$\pm$0.003\\
JPEG 70 & 0.866$\pm$0.005 & 0.862$\pm$0.004\\
JPEG 50 & 0.835$\pm$0.005 & 0.823$\pm$0.004\\
JPEG 30 & 0.771$\pm$0.002 & 0.746$\pm$0.003\\
down $\times$0.75 & 0.840$\pm$0.004 & 0.821$\pm$0.006\\
down $\times$0.5 & 0.835$\pm$0.001 & 0.818$\pm$0.005\\
blur $\sigma{=}1$ & 0.816$\pm$0.001 & 0.805$\pm$0.020\\
\hline\end{tabular}
\end{table}


\paragraph{Results.} Table~\ref{tab:main}. Closing the shortcut raises mean
AUROC (equalized versus as stored), showing that the shortcut otherwise costs
cross-generator accuracy. The detector is near-perfect on the Stable
Diffusion family and weaker on GAN- and ADM-based generators. Accuracy at the
fixed 0.5 cut on the pooled test is only 0.55--0.59 although ranking is good,
the calibration-under-shift problem of \S\ref{sec:fpr}. Robustness degrades
with JPEG quality (Table~\ref{tab:robust}).

\paragraph{Accurate tier.} Adding a frozen CLIP ViT-B~\cite{radford2021clip}
branch (three seeds) raises the clean mean AUROC only slightly but makes the
detector markedly more robust to compression, resizing and blur, by more
than the seed-to-seed spread in every condition (Table~\ref{tab:robust}); it
is stronger on Midjourney, GLIDE and BigGAN and weaker on ADM and VQDM
(Table~\ref{tab:main}).

\paragraph{Newer generators.} On a fixed sample of AIGIBench~\cite{li2025aigibench}
(500 real and 500 generated images per set: FLUX.1-dev, SD 3, Midjourney v6,
DALL-E 3, Imagen 3, SDXL, and two sets collected in the wild), both tiers
fall far below their GenImage results, and the fast tier is near chance on
Midjourney v6 and on images collected from social media
(Table~\ref{tab:aigibench}). The accurate tier is better on every set, but
both remain well below Community Forensics. A detector trained on one 2022
generator does not transfer to current ones; this is the setting in which a
threshold set on real photographs, rather than a fixed cut, matters most.

% Generated by experiments/phase4/aigibench_table.py -- do not edit by hand.
\begin{table}[t]\centering\small
\caption{AUROC on AIGIBench 2024--25 generators and in-the-wild sets (500 real + 500 generated each). fakeradar trained on SD v1.4 only, compression equalized; Community Forensics re-evaluated as stored.}
\label{tab:aigibench}
\begin{tabular}{l c c c}\hline
Test set & fast & accurate & CommFor\\\hline
FLUX.1-dev & 0.722 & 0.821 & 0.974\\
SD 3 & 0.788 & 0.822 & 0.976\\
Midjourney v6 & 0.494 & 0.666 & 0.994\\
DALL-E 3 & 0.859 & 0.898 & 0.980\\
Imagen 3 & 0.637 & 0.680 & 0.925\\
SDXL & 0.756 & 0.813 & 0.998\\
CommunityAI (wild) & 0.584 & 0.796 & 0.936\\
SocialRF (wild) & 0.479 & 0.531 & 0.935\\
\hline\textbf{Mean} & 0.665 & 0.753 & 0.965\\
\hline\end{tabular}
\end{table}


\paragraph{Coherence ablation.} Removing the coherence channel changes clean
mean AUROC within seed noise, but costs 0.01--0.03 AUROC under strong JPEG
and downscaling in both seed pairs (Table~\ref{tab:robust}). We claim no more
than that; similar coherence features appear in public code.

\paragraph{Reference detector.} The public Community Forensics
model~\cite{park2025community}, scored on the same test images without
retraining, reaches a clearly higher mean AUROC and stays stronger under
every perturbation (Tables~\ref{tab:main} and~\ref{tab:robust}). It was
trained on thousands of generators, very likely including GenImage's, so this
is not a same-data comparison; it shows that broad training data, not
architecture, dominates cross-generator accuracy at this scale, and that our
small detector is not competitive on accuracy. Its contribution here is as an
openly measured, CPU-class baseline and as the testbed for the guarantee.

\section{Limitations and Ethics}
We did not retrain other detectors (CNNDetection, NPR, UnivFD) on the same
Stable Diffusion v1.4 data, so our accuracy comparison is against a
re-evaluated public model trained on far more generators, not a same-data
comparison; it bounds how far behind a small single-generator detector is,
but does not rank architectures.
Training and the main benchmark use 2022--2023 generators; the AIGIBench
check below is a single sample of 2024--25 generators, not a full evaluation. The guarantee assumes the real photographs being judged
resemble the calibration photographs of their stratum; it must be
recalibrated when cameras or platforms change, and it bounds false
accusations, not misses. Scores are probabilistic signals and must not be the
sole evidence about a person. Detection research is dual-use; we publish
weights, raw scores and the negative result so that claims can be checked.

\section{Conclusion}
Single-image lens physics was not usable as evidence with today's monocular
estimators. A guarantee on false accusations, however, needs no knowledge of
generators: real-only, platform-stratified conformal thresholds held on unseen
phones and platforms. We suggest reporting the realised false-accusation rate
on held-out devices alongside AUROC.

\paragraph{Reproducibility.} Code, configurations, the Gate M and Gate B
reports, raw per-image scores, and weights
(\texttt{huggingface.co/JitendraJha98/fakeradar-fast}, CC-BY-NC-SA-4.0) are
released.

\bibliographystyle{plain}
\bibliography{references}

\end{document}
